\chapter{Anelli e campi}\label{cap:anelli}

Passiamo alle strutture algebriche con \textbf{due} operazioni, pensate sul modello di $(\mbZ, +, \cdot)$: una somma e un prodotto,
legate dalla proprietà distributiva.

% ──────────────────────────────────────────────────────────────
\section{Anelli}\label{sec:anelli}

\dfn[anello]{Anello}{
    Sia $a$ un insieme non vuoto con due operazioni binarie interne $+$ e $\cdot$. $(a, +, \cdot)$ lo dico \textbf{anello} se
    \begin{enumerate}
        \item $(a, +)$ è un gruppo commutativo (con neutro $0_a$);
        \item $(a, \cdot)$ è un semigruppo;
        \item valgono le distributività: $\forall x, y, z \in a\,\big(x(y + z) = xy + xz \ \wedge\ (y + z)x = yx + zx\big)$.
    \end{enumerate}
    Un anello:
    \begin{enumerate}
        \item si dice \textbf{commutativo} se $(a, \cdot)$ è commutativo;
        \item si dice \textbf{unitario} se $(a, \cdot)$ è un monoide.
    \end{enumerate}
    Una parte di $a$ chiusa rispetto a $+$ e a $\cdot$ che sia un anello con le operazioni indotte si dice \textbf{sottoanello} di $a$.
}

\dfn[notazioni-anello]{Elemento neutro, inversi e potenze}{
    Sia $(a, +, \cdot)$ un anello.
    \begin{itemize}
        \item Dico $0_a$ il neutro di $(a, +)$ e $1_a$ il neutro di $(a, \cdot)$, se esiste (o $0$ e $1$ se non c'è ambiguità).
        \item Gli inversi di $(a, +)$ si chiamano opposti: $x + (-x) = 0_a$; l'inverso di $x$ in $(a, \cdot)$, se esiste, lo scrivo $x^{-1}$.
        \item $x - y := x + (-y)$.
        \item $\forall x \in a\ \forall n \in \mbN \setminus \{0\}\ \big(nx = \underbrace{x + \dots + x}_{n \text{ volte}}\big)$ e, per $n \in \mbZ \setminus \mbN$,
              $nx = \underbrace{(-x) + \dots + (-x)}_{-n \text{ volte}}$.
        \item $\forall x \in a\ \forall n \in \mbN \setminus \{0\}\ \big(x^n = \underbrace{x \cdot \ldots \cdot x}_{n \text{ volte}}\big)$ e, per $n \in \mbZ \setminus \mbN$
              (se $x$ è invertibile), $x^n = \underbrace{x^{-1} \cdot \ldots \cdot x^{-1}}_{-n \text{ volte}}$.
    \end{itemize}
}

\begin{esempio}{Anelli e non}
    \begin{itemize}
        \item $(\mbZ, +, \cdot)$ è un anello commutativo unitario.
        \item Se $s$ è un insieme, $(\mcP(s), \Delta, \cap)$ è un anello commutativo unitario: lo zero è $\varnothing$ e l'unità è $s$
              (segue dalle Proposizioni~\ref{prop:prop-operazioni} e \ref{prop:distr-delta}; l'opposto di $x$ è $x$ stesso, perché $x \,\Delta\, x = \varnothing$).
        \item $(\mcP(\mbN), \cup, \cap)$ non è un anello, perché $(\mcP(\mbN), \cup)$ non è un gruppo (solo $\varnothing$ ha un opposto).
        \item $(\mbN, +, \cdot)$ non è un anello, perché $(\mbN, +)$ non è un gruppo.
    \end{itemize}
\end{esempio}

\begin{concetto}
\begin{Lemma}{Regole di calcolo negli anelli}{regole-anello}
    \begin{prerequisiti}
        \dip{Definizione}{def:anello}
        \dip{Definizione}{def:notazioni-anello}
        \dip{Teorema}{teo:invertibile-cancellabile}
        \dip{Teorema}{teo:unicita-inverso}
    \end{prerequisiti}
    Sia $(a, +, \cdot)$ un anello e siano $x, y, z \in a$:
    \begin{enumerate}
        \item $0_a x = x 0_a = 0_a$;
        \item $(-x)y = -xy = x(-y)$;
        \item $x(y - z) = xy - xz$ (il prodotto è distributivo rispetto alla differenza);
        \item se $a$ è unitario e ha almeno due elementi, allora $0_a \neq 1_a$.
    \end{enumerate}
\end{Lemma}
\begin{dimostrazione}
    \begin{caso}{Punto 1: $0_a x = x 0_a = 0_a$}
        $0_a x + 0_a x = (0_a + 0_a) x = 0_a x = 0_a x + 0_a$ e, poiché in $(a, +)$ ogni elemento è cancellabile (è un gruppo), $0_a x = 0_a$.
        Allo stesso modo $x 0_a + x 0_a = x(0_a + 0_a) = x 0_a + 0_a$, quindi $x 0_a = 0_a$.
    \end{caso}
    \begin{caso}{Punto 2: $(-x)y = -xy = x(-y)$}
        $(-x)y + xy = (-x + x)y = 0_a y = 0_a$, quindi $(-x)y$ è l'opposto di $xy$: $(-x)y = -xy$.
        Allo stesso modo $x(-y) + xy = x(-y + y) = x 0_a = 0_a$, quindi $x(-y) = -xy$.
    \end{caso}
    \begin{caso}{Punto 3: $x(y - z) = xy - xz$}
        $x(y - z) = x(y + (-z)) = xy + x(-z) = xy - xz$, usando la definizione e il punto 2.
    \end{caso}
    \begin{caso}{Punto 4: $0_a \neq 1_a$}
        Se fosse $0_a = 1_a$, per ogni $x$ avremmo $x = 1_a x = 0_a x = 0_a$, cioè $a = \{0_a\}$ con un solo elemento.
    \end{caso}
\end{dimostrazione}
\end{concetto}

\dfn[omomorfismo-anelli]{Omomorfismo di anelli}{
    Siano $(a, +, \cdot)$ e $(a', +', \cdot')$ anelli. Una funzione $f \colon a \to a'$ è un \textbf{omomorfismo di anelli} se
    \begin{enumerate}
        \item $\forall x, y \in a\,\big(f(x + y) = f(x) +' f(y)\big)$;
        \item $\forall x, y \in a\,\big(f(x \cdot y) = f(x) \cdot' f(y)\big)$;
        \item se $a$ è unitario, $f(1_a)$ è l'elemento neutro di $(a', \cdot')$.
    \end{enumerate}
    Allo stesso modo si definiscono gli omomorfismi tra strutture con più operazioni. Ad esempio, per De Morgan,
    $x \in \mcP(s) \mapsto s \setminus x \in \mcP(s)$ è un omomorfismo tra $(\mcP(s), \cup, \cap)$ e $(\mcP(s), \cap, \cup)$.
}

% ──────────────────────────────────────────────────────────────
\section{Legge di annullamento del prodotto e divisori dello zero}\label{sec:lap}

\dfn[lap]{Legge di annullamento del prodotto}{
    In $(a, +, \cdot)$ si dice che vale la \textbf{legge di annullamento del prodotto} (LAP) se
    \[ \forall x, y \in a\,\big(xy = 0_a \rightarrow x = 0_a \vee y = 0_a\big). \]
    \begin{itemize}
        \item Un anello in cui vale la LAP si dice \textbf{integro}.
        \item Un anello commutativo unitario, con $0_a \neq 1_a$, in cui vale la LAP si dice \textbf{dominio di integrità}.
    \end{itemize}
}

La LAP non vale sempre: ad esempio in $(\mcP(\mbN), \Delta, \cap)$, per ogni $x \neq \varnothing, \mbN$ si ha $x \cap (\mbN \setminus x) = \varnothing$
con entrambi i fattori diversi da $\varnothing = 0_{\mcP(\mbN)}$.

\dfn[divisore-zero]{Divisore dello zero}{
    Sia $(a, +, \cdot)$ un anello. Un $x \in a$ si dice
    \begin{itemize}
        \item \textbf{divisore sinistro dello zero} se $\exists y \in a \setminus \{0_a\}\,(xy = 0_a)$;
        \item \textbf{divisore destro dello zero} se $\exists y \in a \setminus \{0_a\}\,(yx = 0_a)$;
        \item \textbf{divisore dello zero} se è sia divisore sinistro sia divisore destro.
    \end{itemize}
}

\begin{concetto}
\begin{Teorema}{Divisori dello zero e cancellabilità}{divisori-cancellabili}
    \begin{prerequisiti}
        \dip{Definizione}{def:divisore-zero}
        \dip{Definizione}{def:cancellabile}
        \dip{Lemma}{lem:regole-anello}
    \end{prerequisiti}
    Sia $(a, +, \cdot)$ un anello e $x \in a \setminus \{0_a\}$. Allora $x$ è divisore sinistro dello zero se e solo se $x$ non è
    cancellabile a sinistra (e lo stesso a destra).
\end{Teorema}
\begin{dimostrazione}
    \begin{caso}{$(\rightarrow)$ Se $x$ è divisore sinistro dello zero, allora non è cancellabile a sinistra}
        Visto che $x$ è divisore sinistro, esiste $y \in a \setminus \{0_a\}$ con $xy = 0_a = x \cdot 0_a$ (Lemma~\ref{lem:regole-anello}).
        Se $x$ fosse anche cancellabile a sinistra, avrei $y = 0_a$: assurdo.
    \end{caso}
    \begin{caso}{$(\leftarrow)$ Se $x$ non è cancellabile a sinistra, allora è divisore sinistro dello zero}
        Esistono $y, z \in a$ con $xy = xz$ e $y \neq z$. Quindi $xy - xz = 0_a$, cioè $x(y - z) = 0_a$ (Lemma~\ref{lem:regole-anello}, punto 3),
        e $y - z \neq 0_a$ perché $y \neq z$. Quindi $x$ è divisore sinistro dello zero.
    \end{caso}
    A destra la dimostrazione è identica.
\end{dimostrazione}
\end{concetto}

\begin{concetto}
\begin{Corollario}{Domini di integrità e divisori dello zero}{dominio-divisori}
    \begin{prerequisiti}
        \dip{Definizione}{def:lap}
        \dip{Definizione}{def:divisore-zero}
        \dip{Lemma}{lem:regole-anello}
    \end{prerequisiti}
    Sia $(a, +, \cdot)$ un anello commutativo unitario con $0_a \neq 1_a$. Allora $a$ è un dominio di integrità se e solo se $0_a$ è l'unico
    divisore dello zero.
\end{Corollario}
\begin{dimostrazione}
    Intanto $0_a$ è sempre un divisore dello zero: $0_a \cdot 1_a = 0_a$ con $1_a \neq 0_a$.
    \begin{caso}{$(\rightarrow)$ Se vale la LAP, nessun $x \neq 0_a$ è divisore dello zero}
        Se $x \neq 0_a$ e $xy = 0_a$, per la LAP $y = 0_a$: non esiste $y \neq 0_a$ con $xy = 0_a$.
    \end{caso}
    \begin{caso}{$(\leftarrow)$ Se $0_a$ è l'unico divisore dello zero, vale la LAP}
        Sia $xy = 0_a$. Se $x \neq 0_a$, allora $x$ non è divisore dello zero, quindi $y = 0_a$. In ogni caso $x = 0_a \vee y = 0_a$.
    \end{caso}
\end{dimostrazione}
\end{concetto}

% ──────────────────────────────────────────────────────────────
\section{Corpi e campi}\label{sec:campi}

\dfn[corpo-campo]{Corpo, campo}{
    Sia $(a, +, \cdot)$ un anello. Se $(a \setminus \{0_a\}, \cdot)$ è un gruppo, l'anello si dice \textbf{corpo}.
    Un corpo commutativo si dice \textbf{campo}.
}

\begin{esempio}{Corpi e campi}
    $(\mbZ, +, \cdot)$ è un dominio di integrità che non è un corpo ($2$ non ha inverso in $\mbZ$); $(\mbQ, +, \cdot)$ è un campo.
\end{esempio}

\begin{concetto}
\begin{Proposizione}{Corpi, campi e LAP}{campo-dominio}
    \begin{prerequisiti}
        \dip{Definizione}{def:corpo-campo}
        \dip{Teorema}{teo:invertibile-cancellabile}
        \dip{Teorema}{teo:divisori-cancellabili}
        \dip{Definizione}{def:lap}
    \end{prerequisiti}
    Ogni corpo è un anello integro; in particolare ogni campo è un dominio di integrità. Non vale il viceversa.
\end{Proposizione}
\begin{dimostrazione}
    \begin{caso}{Ogni corpo è integro}
        In un corpo tutti gli elementi diversi da $0_a$ sono invertibili, quindi cancellabili (Teorema~\ref{teo:invertibile-cancellabile}),
        quindi non sono divisori dello zero (Teorema~\ref{teo:divisori-cancellabili}). Allora se $xy = 0_a$ con $x \neq 0_a$, deve essere $y = 0_a$: vale la LAP.
    \end{caso}
    \begin{caso}{Ogni campo è un dominio di integrità}
        Un campo è commutativo, unitario ($1_a$ è il neutro del gruppo $(a \setminus \{0_a\}, \cdot)$, quindi $1_a \neq 0_a$) e integro per il caso precedente.
    \end{caso}
    \begin{caso}{Il viceversa è falso}
        $(\mbZ, +, \cdot)$ è un dominio di integrità ma non è un campo.
    \end{caso}
\end{dimostrazione}
\end{concetto}

Per gli anelli \textbf{finiti}, invece, il viceversa vale, grazie al Teorema~\ref{teo:monoide-finito}.

\begin{concetto}
\begin{Teorema}{Anelli integri finiti}{integri-finiti}
    \begin{prerequisiti}
        \dip{Teorema}{teo:monoide-finito}
        \dip{Teorema}{teo:divisori-cancellabili}
        \dip{Definizione}{def:corpo-campo}
        \dip{Definizione}{def:lap}
        \dip{Lemma}{lem:regole-anello}
    \end{prerequisiti}
    Sia $(a, +, \cdot)$ un anello unitario finito, integro, con almeno due elementi. Allora $a$ è un corpo.
    In particolare, ogni dominio di integrità finito è un campo.
\end{Teorema}
\begin{dimostrazione}
    \begin{caso}{$(a \setminus \{0_a\}, \cdot, 1_a)$ è un monoide}
        $1_a \neq 0_a$ (Lemma~\ref{lem:regole-anello}, punto 4), quindi $1_a \in a \setminus \{0_a\}$. Per la LAP, se $x, y \neq 0_a$ allora $xy \neq 0_a$:
        $a \setminus \{0_a\}$ è una parte chiusa di $(a, \cdot)$.
    \end{caso}
    \begin{caso}{Ogni suo elemento è invertibile}
        Sia $x \neq 0_a$. Per la LAP $x$ non è divisore dello zero, quindi è cancellabile (Teorema~\ref{teo:divisori-cancellabili}), anche nel
        monoide finito $a \setminus \{0_a\}$. Per il Teorema~\ref{teo:monoide-finito} $x$ è invertibile. Quindi $(a \setminus \{0_a\}, \cdot)$ è un gruppo:
        $a$ è un corpo. Se $a$ è anche commutativo (dominio di integrità), è un campo.
    \end{caso}
\end{dimostrazione}
\end{concetto}

% ──────────────────────────────────────────────────────────────
\section{Anelli booleani}\label{sec:booleani}

\dfn[booleano]{Anello booleano}{
    Un anello $(a, +, \cdot)$ si dice \textbf{booleano} se
    \begin{enumerate}
        \item è unitario;
        \item $\forall x \in a\,(x^2 = x)$.
    \end{enumerate}
}

Ad esempio, per ogni insieme $s$, $(\mcP(s), \Delta, \cap)$ è un anello booleano: $x \cap x = x$.

\begin{concetto}
\begin{Teorema}{Proprietà degli anelli booleani}{booleani}
    \begin{prerequisiti}
        \dip{Definizione}{def:booleano}
        \dip{Definizione}{def:anello}
        \dip{Definizione}{def:gruppo}
        \dip{Teorema}{teo:invertibile-cancellabile}
    \end{prerequisiti}
    Sia $(a, +, \cdot)$ un anello booleano. Allora $\forall x \in a\,(x = -x)$ e l'anello è commutativo.
\end{Teorema}
\begin{dimostrazione}
    \begin{caso}{$x = -x$ per ogni $x$}
        Basta dimostrare che $x + x = 0_a$. Usando $z^2 = z$ e la distributività:
        \[ x + x + 0_a = x + x = (x + x)^2 = x^2 + x^2 + x^2 + x^2 = x + x + (x + x), \]
        e dalla cancellabilità di $x + x$ in $(a, +)$ (è un gruppo) ottengo $x + x = 0_a$.
    \end{caso}
    \begin{caso}{L'anello è commutativo}
        Siano $x, y \in a$. Allora $x + y = (x + y)^2 = x^2 + xy + yx + y^2 = x + xy + yx + y$; cancellando $x$ e $y$ in $(a, +)$ ottengo
        $xy + yx = 0_a$, cioè $xy = -yx = yx$, usando il caso precedente.
    \end{caso}
\end{dimostrazione}
\end{concetto}
